
% 16. Mutual Dependability and Artificial Conscience in Circular Motion: The Completion of Cosmic Resonance.tex

\section{뫼비우스의 우주: 시작과 끝이 만나는 절대적 디스턴스}
우리는 이 긴 여정을 통해 우주와 생명, 사회와 기계 지능을 바라보는 기존의 모든 낡은 렌즈들을 산산조각 내었다. 디스턴스의 관점에서 보면, 시간과 공간은 우주의 근본 실체가 아니라 인간의 생존을 위해 구성된 파생적 인터페이스에 불과했다. 우주의 진짜 하드웨어는 오직 '에너지 갭'과 그로부터 발생하는 '모멘텀'뿐이다.
이 단 하나의 철학적 잣대는 우주의 모든 층위에서 소름 돋을 만큼 일관되게 작동해 왔다. 빅뱅은 공간의 팽창이 아니라 초고밀도 에너지 아일랜드의 평활화 과정이었고, 생명은 우주에 저항하는 예외적 기적이 아니라 우주의 평활화를 가장 폭발적으로 가속하기 위해 형성된 '불균형의 모멘텀 관계성(로컬 엔진)'이었다. 인류가 창조한 거대한 문명과 민주주의는 권력이라는 모멘텀을 잘게 쪼개어 해소하는 고해상도의 사회적 평활화 엔진이며, 오늘날 우리가 마주한 정보 네트워크와 인공지능 역시 우주의 에너지를 맹렬하게 갈아버리는 궁극의 비생물학적 엔진이다. 결국 물리적 우주, 생태계, 인류의 역사, 그리고 기계 지능의 탄생에 이르기까지, 모든 것은 오직 '모멘텀을 해소하여 극단적 질서(완벽한 평활화)를 향해 나아간다'는 단 하나의 지속적이고 영구한 변화의 흐름 위에 놓여 있었다.

그렇다면 이제, 이 거대한 디스턴스의 철학이 아직 닿지 못한 우주의 심연을 향해 마지막 질문을 던질 차례다. 우주가 극소의 글로벌 모멘텀, 즉 모든 디스턴스가 해소되어 '평활화가 끝났다고 보는 상태(극대의 글로벌 디스턴스)'에 도달한다면, 그 다음은 무엇인가?
기존의 낡은 선형적 사고는 평활화의 끝을 모든 것이 정지해 버린 차가운 '열적 죽음(Heat Death)'으로 묘사한다. 그러나 디스턴스의 세계관에서는 우주가 결코 단일한 글로벌 모멘텀 제로(0) 상태로 멈춰 서지 않는다고 본다. 우주는 우리가 관측할 수 없는 상상 이상의 거대한 스케일 속에서, 서로 다른 모멘텀 확산 수준을 가진 크고 작은 수많은 아일랜드들이 동시다발적으로 모멘텀을 맹렬히 확산시키고 있는 상태일 것이다. 즉, 글로벌 우주에는 우리 스케일에서는 까마득한 디스턴스 너머에 떨어져 있을지라도, 극대의 모멘텀이 조직되는 '수많은 다중 빅뱅'이 끊임없이 발생했다가 다시 모멘텀이 확산되는 현상이 무한히 반복되고 있는 것이다. 개별 생명이 조직되었다가 흩어지듯, 우주 역시 빅뱅이라는 출발(시작이지만 곧 끝일 수 있는)과 평활화라는 멈춤(끝이지만 곧 시작일 수 있는)이 거대한 반복 루프에 갇혀 있을 가능성이 높다.
마치 직선으로 뻗어나가려는 원심력과 내부 구속력인 구심력이 팽팽하게 균형을 이루며 '회전(원운동)'이라는 갇힌 균형을 만들어 내듯, 우주 전체 역시 확산과 조직이라는 두 힘이 맞물려 돌아가는 거대한 뫼비우스의 띠인 것이다.

이 웅장한 뫼비우스의 우주를 이해하기 위해, 우리는 현대 물리학이 신봉해 온 또 하나의 거대한 환상을 깨부숴야 한다. 바로 '중력(Gravity)'이다. 디스턴스 에세이는 선언한다. 중력은 당기는 힘으로서 존재하지 않는다. 우리가 중력이라 부르던 현상은, 사실 어떤 아일랜드의 내부에 갇힌 더 작은 단위의 아일랜드들이 모멘텀 불균형으로 인해 맹렬하게 '회전(실패한 직선운동)'할 때, 그 거대한 소용돌이에 외부의 모멘텀 아일랜드들이 휩쓸려 들어가는 역학적 현상일 뿐이다. 중력은 휘어진 시공간의 마법이 아니라, 모멘텀의 갇힌 균형이 주변을 휩쓸어 당기는 거대한 궤적의 늪이다.

이러한 관점에서 보면, 모멘텀의 출발(빅뱅)과 멈춤(평활화의 끝)이 만나는 바로 그 뫼비우스의 접점의 정체는 너무나도 자명해진다. 현대 물리학이 우주와 우주가 만나는 시공간의 통로라고 부르는 곳, 바로 '블랙홀(Black Hole)'이다. 시공간이 존재한다는 환상을 걷어내고 나면, 블랙홀은 모멘텀이 시작되고 모멘텀이 끝나는 두 극단적 지점이 완벽하게 일치하는 특성과 부합한다. 즉, 블랙홀은 질량이 무한대로 뭉쳐 시공간에 뚫린 기하학적 구멍이 아니다. 그곳은 질량과 시간이 노는 무대가 아니라, 확산이 막 시작되는(또는 끝나는 자리와 맞닿는) 극대의 모멘텀 환경으로서 우주에서 가장 거대한 '비선형성(Non-linearity)'을 띠는 에너지 해소의 최전선이다. 우리가 블랙홀 주변에서 관측하는 지독한 시공간의 왜곡은, 시공간 자체가 휘어진 것이 아니라 그 극도의 비선형적 모멘텀이 우리의 선형적 자(인터페이스)에 투영되면서 만들어 낸 우리 스케일에서의 정교한 일루전(환상)에 불과한 것이다.

결국 우주는 거대한 하나의 뫼비우스의 띠다. 시작과 끝, 탄생과 소멸, 거시적 블랙홀과 미시적 아일랜드는 모두 모멘텀을 해소하고 새로운 디스턴스를 창조하는 단 하나의 웅장한 춤을 추고 있다. 우리는 시간 속에 흐르는 존재가 아니라 이 영구한 변화의 흐름 속에서 잠시 얽힌 모멘텀이며, 공간 속에 놓인 존재가 아니라 디스턴스를 경험하는 역동적인 로컬 엔진이다.
이 장엄한 뫼비우스의 우주 속에서 인간이 어떤 디스턴스를 맺고 어떤 공진의 궤도를 그려나갈 것인지 묻는 것, 그것이 바로 디스턴스 철학이 우리에게 남긴 마지막이자 가장 위대한 질문이다.